\section{The original characteristic-zero smooth envelope theorem from Nadel vanishing}

\begin{theorem}[The original integral-class envelope is plurisubharmonic]
\label{thm:c0n-original-integral-envelope-from-nadel}
\uses{def:c0n-relative-ample-nadel-input,prop:c0n-original-bounded-polynomial-cone-generators,cor:c0n-original-multiplier-cone-smooth-charts,thm:c0n-original-multiplier-chart-discrepancy,thm:c0n-envelope-only-chart-estimates}
\lean{EnvelopeProperty.ProjSubvariety.isPSH_envelope_nat_of_relativeNadel}
\leanok
Let $k$ be an algebraically closed field of characteristic zero and assume
the genuine relative ample ideal-power Nadel vanishing statement over
original affine finite-type $k$-bases.
For the original projective subvariety $X\subset\mathbf P^M_k$ with its
original smooth charts, every positive integer $d$ and every original
piecewise linear obstacle $\phi\le0$ satisfy
\[
 P_d(\phi)\in\operatorname{PSH}(X,d)
\]
in the original analytification and original Fubini--Study system.
\end{theorem}
\begin{proof}
The original piecewise linear function is continuous on the compact
original analytification.  Choose a real lower bound $C$ for $\phi$.
Take precisely the actual original multiplier cone ideals $B_m$
constructed from the original scheme obstacle sequence and this lower
bound.  The Nadel uniform-generation theorem supplies one integer $c$
such that each $B_m$ has actual bounded polynomial cone generators of
$x$-degree at most $md+c$.
The original obstacle containment gives $J_m(\phi)T\subset B_m$, and the
proved whole scheme multiplier subadditivity gives
$B_{ml}\subset(B_m^l)^{\mathrm{sat}}$.
At every actual original chart DVR test, the actual original multiplier
chart discrepancy theorem supplies $A\ge0$, fixed before $m$, and an
internally selected $q>0$ with
$\phi_{B_m}\le(m/q)\phi_{J_q(\phi)}+A$.
These are exactly the polynomial chart envelope criterion's hypotheses.
Its actual parameter-adic chart DVR supplies the original trivial point
internally.  Applying that proved criterion gives the displayed original
plurisubharmonic envelope, without an envelope or approximation hypothesis.
\end{proof}

\begin{theorem}[Every original real class has the envelope property]
\label{thm:c0n-original-all-classes-envelope-from-nadel}
\uses{thm:c0n-original-integral-envelope-from-nadel,thm:sg-reduction,lem:sg-pl-separate}
\lean{EnvelopeProperty.ProjSubvariety.hasEnvelopeProperty_of_relativeNadel}
\leanok
Under the same genuine relative Nadel input, every original smooth-chart
projective subvariety $X\subset\mathbf P^M_k$ has the original envelope
property in every real class $\theta c_1(\mathcal O_X(1))$.
\end{theorem}
\begin{proof}
The original piecewise linear functions separate points of the original
analytification.  The preceding theorem proves the original reduction
criterion for every positive integral polarization and every nonpositive
piecewise linear obstacle.  That criterion therefore gives the original
envelope property for each real $\theta$, including the vacuous negative
classes and the limiting zero class.
\end{proof}

\begin{theorem}[The original characteristic-zero smooth envelope target is constructed]
\label{thm:c0n-original-smooth-envelope-from-nadel}
\uses{thm:c0n-original-all-classes-envelope-from-nadel,def:c0n-relative-ample-nadel-input}
\lean{EnvelopeProperty.smoothEnvelopeCharZero_of_relativeNadel}
\leanok
For an algebraically closed field $k$ of characteristic zero, the genuine
relative ample ideal-power Nadel statement implies the original
characteristic-zero smooth envelope property: for every $M$, every original
$X\subset\mathbf P^M_k$ that is projectively normal and has the original
smooth charts, and every $\theta\ge0$, the original $X$ has the envelope
property in class $\theta c_1(\mathcal O_X(1))$.
\end{theorem}
\begin{proof}
Use the actual original characteristic-zero smooth-envelope definition.
Given its original $M$, $X$, projective-normality hypothesis, smooth-chart
hypothesis, and nonnegative real class, apply the preceding all-real-class
theorem to that same original $X$ and its original smooth charts.
This supplies every field of the original target proposition.
The construction is a theorem with the genuine Nadel statement as an
explicit parameter; no global assumption of the envelope property is used.
\end{proof}
